\subsection{Caratteristiche}
Innanzitutto, è stato necessario definire le modalità di gestione delle due condizioni sperimentali del sistema (la \textit{Versione A} originale e la \textit{Versione B} riprogettata). Per confrontare l'usabilità dei due sistemi, si è optato per un approccio \textit{within-subjects}, prevedendo che ogni partecipante provasse entrambe le condizioni. 

La scelta di questa metodologia rispetto all'altro approccio \textit{between-subjects}, che avrebbe escluso a monte il problema dell'effetto dell'apprendimento, è stata motivata dal fatto che la riprogettazione ha comportato una ristrutturazione in parte profonda, introducendo cambiamenti sostanziali anche a livello funzionale in alcune parti del sistema. Di conseguenza, è stato ritenuto più significativo valutare come lo stesso utente si ponesse di fronte a logiche di navigazione differenti. Oltre a questo aspetto, l'approccio scelto ha offerto il vantaggio di ottimizzare le risorse, richiedendo un numero più contenuto di partecipanti. 

Poiché l'apprendimento intrinseco nel passaggio da una versione all'altra avrebbe potuto falsare i dati raccolti, il problema è stato controllato applicando la tecnica del \textit{counterbalancing}: il campione di utenti è stato infatti diviso approssimativamente a metà, per cui un gruppo ha testato prima la condizione sperimentale A e successivamente la B, mentre il secondo gruppo ha seguito l'ordine inverso. 